\chapter{Construction analytique des fonctionnelles de la tour triadique}
\label{chap:funcionales-espectrales-triadicos}

La synthèse précédente a situé les constantes spectrales au sein de la
tour triadique. Ce bloc contient leur construction analytique primaire :
suite d'opérateurs autoadjoints sur des cycles d'ordre \(3^m\), limite de
traces de chaleur et transformée de Mellin. Les constantes classiques apparaissent
comme évaluations, parties finies ou dérivées de la fonctionnelle méromorphe
résultante, non comme des préfixes reconnus dans un catalogue.

La normalisation de l'opérateur $D_m$ contient explicitement $\pi$ dans l'échelle
qui relie le laplacien discret au noyau de chaleur gaussien. Cette
sous-section constitue donc une extension analytique de la branche
archimédienne déjà construite. Les constantes spectrales ultérieures s'obtiennent
simultanément à partir de la fonctionnelle résultante, avec cette normalisation pour hypothèse
de départ.

\section{Limites inductives et échelle de raffinement}

Le prolongement spectral exige de distinguer deux niveaux analytiques. Le
premier correspond à des familles bornées et compatibles ; le second, à des
opérateurs non bornés, pour lesquels l'entrelacement algébrique ne
résout pas à lui seul les problèmes de domaine.

\begin{theorem}[Limite inductive bornée]
\label{pf84:E03-06}
Soient \(J_m:\mathcal H_m\to\mathcal H_{m+1}\) des isométries et
\(A_m\in\mathcal B(\mathcal H_m)\) des opérateurs autoadjoints tels que
\[
 A_{m+1}J_m=J_mA_m,
 \qquad
 \sup_m\lVert A_m\rVert<\infty.
\]
Il existe un unique opérateur borné autoadjoint \(A\) sur
\(\mathcal H_\infty=\varinjlim(\mathcal H_m,J_m)\) dont la restriction à chaque
niveau est \(A_m\).
\end{theorem}

\begin{proof}
Sur l'union algébrique dense, on définit
\[
 A[J_{m,\infty}\xi]:=J_{m,\infty}A_m\xi.
\]
L'entrelacement rend la définition indépendante du représentant. La
borne uniforme permet de prolonger \(A\) par continuité à
\(\mathcal H_\infty\) ; la symétrie sur le sous-espace dense et le caractère borné
impliquent son autoadjonction.
\end{proof}

\begin{remark}[Contrôle pour les opérateurs non bornés]
\label{pf84:E03-07}
Si les \(A_m\) ne sont pas bornés, l'identité
\(A_{m+1}J_m=J_mA_m\) ne définit qu'une action sur une union de domaines.
Pour obtenir un opérateur fermé ou autoadjoint, il faut déclarer les domaines
transportés, prouver la fermabilité et contrôler les indices de
déficience, ou bien démontrer l'autoadjonction essentielle sur un cœur commun
compatible. Omettre ces obligations constitue une erreur de type, non une
abréviation de la preuve.
\end{remark}

Soit maintenant \(\mathcal S_\infty\) un ensemble non vide de sections
survivantes compatibles et soit \(\mathcal P_m\) l'ensemble fini de leurs
préfixes de longueur \(m\). Nous écrivons
\[
 \mathcal H_m:=\ell^2(\mathcal P_m),
 \qquad
 E_m(\gamma,\gamma')=
 \begin{cases}
 1,&\gamma'\text{ prolonge }\gamma,\\
 0,&\text{sinon},
 \end{cases}
\]
et
\[
 d_m(\gamma):=
 \#\{\gamma'\in\mathcal P_{m+1}:\gamma'\mapsto\gamma\},
 \qquad
 D_m^{\deg}:=E_mE_m^*=\operatorname{diag}(d_m(\gamma)).
\]
Nous nous restreignons au sous-espace vivant \(d_m(\gamma)>0\), et chaque préfixe de
niveau \(m+1\) possède un parent unique.

\begin{theorem}[Isométrie normalisée de prolongement]
\label{pf84:C01}
L'opérateur
\[
 J_m:=E_m^*(D_m^{\deg})^{-1/2}:
 \mathcal H_m\longrightarrow\mathcal H_{m+1}
\]
est une isométrie, et
\[
 J_me_\gamma=
 \frac1{\sqrt{d_m(\gamma)}}
 \sum_{\gamma'\mapsto\gamma}e_{\gamma'}.
\]
\end{theorem}

\begin{proof}
L'unicité du parent donne
\[
 J_m^*J_m
 =(D_m^{\deg})^{-1/2}E_mE_m^*(D_m^{\deg})^{-1/2}=I.
\]
La formule sur \(e_\gamma\) s'obtient en développant la colonne
correspondante de \(E_m^*\).
\end{proof}

En posant
\[
 \mathcal K_m:=\mathcal H_{m+1}\ominus J_m\mathcal H_m,
 \qquad
 \mathcal H_{\mathrm{surv}}
 \cong\mathcal H_0\oplus\bigoplus_{m\ge0}\mathcal K_m,
\]
on définit l'opérateur de raffinement
\[
 \mathcal R|_{\mathcal H_0}=0,
 \qquad
 \mathcal R|_{\mathcal K_m}=3^{m+1}I_{\mathcal K_m},
\]
de domaine maximal
\[
 \Dom(\mathcal R)=
 \mathcal H_0\oplus
 \left\{(\xi_m)_{m\ge0}:
 \sum_{m\ge0}3^{2m+2}\lVert\xi_m\rVert^2<\infty\right\}.
\]

\begin{theorem}[Opérateur de raffinement]
\label{pf84:C02}
Si chaque \(\mathcal K_m\) est de dimension finie, \(\mathcal R\) est positif,
autoadjoint et à résolvante compacte. Son spectre mesure la profondeur de
raffinement ; il ne constitue pas un spectre de nombres premiers ni de zéros.
\end{theorem}

\begin{proof}
La décomposition orthogonale réduit \(\mathcal R\) à des blocs scalaires réels
sur son domaine maximal, ce qui prouve la positivité et l'autoadjonction. Les
valeurs propres \(3^{m+1}\) s'échappent à l'infini et leurs multiplicités sont finies ;
la résolvante est donc compacte. Un bloc de dimension infinie à
valeur propre bornée réfuterait cette dernière conclusion.
\end{proof}

\section{Opérateurs entier, premier et orbital}

Soit \(\mathscr W_{\mathrm{bal}}\) l'ensemble des mots ternaires
équilibrés réduit au moyen du théorème
\cref{pf84:C03}. Sur \(\ell^2(\mathscr W_{\mathrm{bal}})\), l'opérateur
\[
 Be_w=b(w)e_w
\]
est pris sur son domaine maximal. Sa restriction au sous-espace des entiers
positifs est notée \(Q\).

\begin{theorem}[Opérateurs entier et positif]
\label{pf84:C04}
L'opérateur \(B\) est autoadjoint, possède un spectre simple \(\ZZ\) et une
résolvante compacte. L'opérateur \(Q\) possède un spectre simple
\(\NN_{\ge1}\) et une résolvante compacte.
\end{theorem}

\begin{proof}
La bijection de \cref{pf84:C03} énumère chaque entier exactement une fois. Un
opérateur diagonal réel sur son domaine maximal est autoadjoint ; comme ses entrées
s'échappent à l'infini hors des ensembles finis, sa résolvante est compacte. La
restriction positive conserve ces propriétés.
\end{proof}

La preuve classique omise de l'inversion thêta, du prolongement et de l'équation
fonctionnelle est fusionnée une seule fois avec la réalisation triadique déjà construite
dans \cref{pf84:C05}. On y fixe littéralement
\(\Theta_B=1+2\Theta_+\) et \(Z=\mathcal Z_3\) ; ce bloc ne
définit donc pas une seconde origine de \(\zeta\).

Soit \(\mathbb P\) l'ensemble des nombres premiers ordinaires,
\(\mathfrak h_{\mathbb P}:=\ell^2(\mathbb P)\) et
\[
 H_{\mathbb P}e_p=(\log p)e_p.
\]
Dans l'espace de Fock bosonique \(\Gamma_s(\mathfrak h_{\mathbb P})\), les occupations
\((k_p)_p\) sont à support fini.

\begin{theorem}[Représentation de Fock de la factorisation]
\label{pf84:C08}
L'application sur les bases
\[
 U_F\lvert(k_p)_p\rangle=e_{\prod_pp^{k_p}}
\]
se prolonge en un opérateur unitaire
\[
 U_F:\Gamma_s(\mathfrak h_{\mathbb P})
 \longrightarrow\ell^2(\NN_{\ge1})
\]
et satisfait
\[
 U_F\,d\Gamma(H_{\mathbb P})\,U_F^*=\log Q.
\]
\end{theorem}

\begin{proof}
La factorisation unique donne une bijection entre les occupations finies et les entiers
positifs, donc une identification unitaire des bases orthonormales.
De plus,
\[
 \sum_pk_p\log p=\log\!\left(\prod_pp^{k_p}\right),
\]
ce qui prouve l'entrelacement. La structure de départ explicite est le
choix des modes premiers et de l'espace de Fock symétrique.
\end{proof}

Sur
\[
 \mathcal H_{\mathrm{orb}}
 :=\ell^2\{(p,k,\varepsilon):
 p\in\mathbb P,\ k\ge1,\ \varepsilon=\pm1\},
\]
nous définissons
\[
 D_{\mathrm{orb}}e_{p,k,\varepsilon}
 =\varepsilon k\log p\,e_{p,k,\varepsilon},
 \qquad
 W_{\log}e_{p,k,\varepsilon}
 =(\log p)e_{p,k,\varepsilon}.
\]

\begin{proposition}[Trace orbitale]
\label{pf84:C09}
L'opérateur \(D_{\mathrm{orb}}\) est autoadjoint et à résolvante compacte. Si
\(f\in C_c(\RR)\) est bornée, alors
\[
 \Tr\!\left(
 W_{\log}e^{-|D_{\mathrm{orb}}|/2}f(D_{\mathrm{orb}})
 \right)
 =
 \sum_{p,k}(\log p)p^{-k/2}
 \bigl[f(k\log p)+f(-k\log p)\bigr].
\]
\end{proposition}

\begin{proof}
La diagonale est réelle et ses longueurs s'échappent de tout compact avec
multiplicité finie. Le support compact de \(f\) fait que le produit
régulé est de rang fini ; sa trace est la somme des entrées diagonales.
L'indice \(p\) parcourt ici les nombres premiers ordinaires. Le résultat réalise le
côté premier classique, mais n'identifie pas \(D_{\mathrm{orb}}\) à un opérateur
dual des zéros.
\end{proof}

\section{Forme complétée de Guinand--Weil}

Nous fixons
\[
 \mathcal D_{\mathrm W}:=C_c^\infty(\RR),
 \qquad
 \widehat\psi(t)=\int_{\RR}\psi(x)e^{-itx}\,\dd x,
 \qquad
 \psi(x)=\frac1{2\pi}\int_{\RR}\widehat\psi(t)e^{itx}\,\dd t,
\]
et, pour \(\psi,\phi\in\mathcal D_{\mathrm W}\),
\[
 h_{\psi,\phi}(z)
 :=\widehat\psi(z)\,
 \overline{\widehat\phi(\overline z)}.
\]
Sur la droite réelle \(h_{\psi,\psi}=|\widehat\psi|^2\).

\begin{definition}[Classe de test]
\label{pf84:E01}
La classe de test est \(\mathcal D_{\mathrm W}\), stable par dérivation,
réflexion et translation. On utilise l'inverse
\[
 \check h(x)=\frac1{2\pi}\int_{\RR}h(t)e^{itx}\,\dd t.
\]
La régularité et le support compact fournissent une décroissance rapide dans les bandes
et un support compact pour \(\check h\), conditions qui garantissent la
convergence des blocs suivants.
\end{definition}

Sea \(\psi_\Gamma=\Gamma'/\Gamma\). Definimos
\[
 \Irr_\#:=
 \left\{\operatorname{ev}_9(w):
 w\in\mathscr W_9^+,\ 
 P_{\Irr,\#}e_w=e_w\right\}.
\]
Par \cref{pf84:C07}, cette image coïncide avec les nombres premiers ordinaires ; elle n'est
utilisée ici que pour typer l'indice du bloc orbital, sans répéter cette
preuve.
\begin{align}
 \mathcal P(h)&:=h(i/2)+h(-i/2),\\
 \mathcal G(h)&:=-\log\pi\,\check h(0)
 +\frac1{2\pi}\int_{\RR}h(t)
 \Re\psi_\Gamma\!\left(\frac14+\frac{it}{2}\right)\dd t,\\
 \mathcal O_\#(\check h)&:=
 \sum_{p\in\Irr_\#}\sum_{k\ge1}
 (\log p)p^{-k/2}
 \bigl[\check h(k\log p)+\check h(-k\log p)\bigr].
\end{align}

\begin{definition}[Fonctionnelle complétée]
\label{pf84:E02}
Pour \(h=h_{\psi,\phi}\), on fixe
\[
 \mathcal W(h):=\mathcal P(h)+\mathcal G(h)
 -\mathcal O_\#(\check h).
\]
Le bloc orbital provient du sélecteur irréductible du
\cref{chap:producto-nativo} ; les blocs polaire et gamma appartiennent à la
complétion classique de zêta. Cette différence de provenance est conservée dans
la notation et dans la preuve.
\end{definition}

\begin{theorem}[Formule explicite de Guinand--Weil]
\label{pf84:E03} Pour \(\psi,\phi\in C_c^\infty(\RR)\) et \(h=h_{\psi,\phi}\),
\[
 \lim_{T\to\infty}
 \sum_{\substack{\rho:\,\zeta(\rho)=0\\|\Im\rho|\le T}}
 h\!\left(\frac{\rho-1/2}{i}\right)
 =\mathcal W(h),
\]
où les zéros non triviaux sont comptés avec multiplicité et la limite est
prise avec la symétrisation classique.
\end{theorem}

\begin{proof}
Paley--Wiener place \(h_{\psi,\phi}\) dans la classe admissible de la formule
explicite. Les pôles de
\(\Lambda_\zeta(s)=\pi^{-s/2}\Gamma(s/2)\zeta(s)\) produisent
\(\mathcal P\), le facteur archimédien produit \(\mathcal G\) et le produit
d'Euler produit la somme sur les puissances premières. L'identification de cette
dernière somme à \(\mathcal O_\#\) s'effectue uniquement au moyen de
\cref{pf84:E04} ; ainsi, la preuve de primalité native n'est pas dupliquée.
\end{proof}

\begin{proposition}[Réindexation du bloc orbital]
\label{pf84:E04}
Si les irréductibles du produit natif coïncident exactement avec les nombres premiers
ordinaires et que les poids
\((\log p)p^{-k/2}\) et les deux orientations \(\pm k\log p\) sont conservés, alors
\[
 \mathcal O_\#(\check h)=\mathcal O(\check h).
\]
\end{proposition}

\begin{proof}
La bijection entre irréductibles natifs et nombres premiers réindexe la somme sans
modifier les indices, les poids ni les arguments. Un irréductible non premier, un nombre premier
absent ou un changement de poids réfuterait l'identité.
\end{proof}

\begin{remark}[Séparation méromorphe--entière]
\label{pf84:E05}
La fonction \(\Lambda_\zeta\) possède des pôles simples en \(0\) et \(1\), qui
engendrent le bloc polaire. En revanche,
\[
 \xi(s)=\frac12s(s-1)\Lambda_\zeta(s)
\]
est entière : les facteurs \(s(s-1)\) annulent précisément ces pôles. On
n'attribue donc pas de pôles en \(0\) ou \(1\) à \(\xi\).
\end{remark}

\begin{definition}[Forme hermitienne et préopérateur]
\label{pf84:E06}
On définit
\[
 \mathcal B_{\mathrm W}(\psi,\phi)
 :=\mathcal W(h_{\psi,\phi}),
 \qquad
 \mathcal V_{\mathrm W}:=
 \mathcal D_{\mathrm W}/\operatorname{Rad}\mathcal B_{\mathrm W},
\]
et, sur le quotient algébrique, la règle candidate
\[
 D_{\mathrm W,0}[\psi]:=[-i\psi'].
\]
On ne présuppose pas que \(\mathcal B_{\mathrm W}\) soit positive ni que
\(\mathcal V_{\mathrm W}\) soit déjà un espace de Hilbert. La règle n'est
bien définie qu'après avoir prouvé l'invariance du radical.
\end{definition}

\begin{theorem}[Symétrie algébrique et invariances]
\label{pf84:E07}
Le radical est invariant sous \(-i\,\dd/\dd x\) ;
\(D_{\mathrm W,0}\) est bien défini et
\(\mathcal B_{\mathrm W}\)-symétrique. Les translations
\(T_a\psi(x)=\psi(x-a)\) préservent la forme et la réflexion
\(J\psi(x)=\psi(-x)\) satisfait
\[
 JD_{\mathrm W,0}J=-D_{\mathrm W,0}.
\]
\end{theorem}

\begin{proof}
La convention de Fourier donne
\[
 h_{-i\psi',\phi}(z)=h_{\psi,-i\phi'}(z).
\]
L'identité vaut avant l'application de l'un quelconque des trois blocs de
\(\mathcal W\), de sorte que la dérivée préserve le radical et que le
préopérateur est symétrique. Les phases opposées des translations s'annulent
dans \(h\) ; la réflexion change le signe de la dérivée et laisse invariants les
blocs polaire, gamma et orbital. Cette symétrie algébrique n'équivaut pas encore
à l'autoadjonction dans un Hilbert positif.
\end{proof}

\section{Géométrie locale de la droite critique}

\begin{lemma}[Coordonnée de Möbius de la droite critique]
\label{pf84:RH-01}
Pour \(s\in\CC\setminus\{0,1\}\), soit
\[
 \tau(s):=i\frac{s-1}{s}.
\]
Avec \(F(s)=1-s\) et
\(C(s)=\overline s\),
\[
 \tau(Fs)=-\frac1{\tau(s)},
 \qquad
 \tau(Cs)=-\overline{\tau(s)},
\]
et
\[
 \tau(Fs)=\tau(Cs)
 \iff |\tau(s)|=1
 \iff \Re s=\frac12.
\]
\end{lemma}

\begin{proof}
Les deux identités s'obtiennent par substitution directe. Leur égalité
équivaut à \(\tau(s)^{-1}=\overline{\tau(s)}\), c'est-à-dire
\(|\tau(s)|=1\) ; développer cette dernière condition donne
\(\Re s=1/2\). Pour les zéros de zêta, les symétries
fonctionnelle et conjuguée sont en outre requises. L'équivalence paramètre la droite ; elle ne place pas à elle
seule les zéros sur celle-ci.
\end{proof}

\begin{lemma}[Signe local de Hermite--Biehler]
\label{pf84:RH-04}
Soit \(A\) une fonction entière réelle ayant un zéro réel simple \(a\). Pour
\(z=a+iy\), \(y>0\) suffisamment petit,
\[
|A(z)-iA'(z)|^2-|A(z)+iA'(z)|^2
 =-4A'(a)^2y+O(y^2)<0.
\]
La convention locale associée à \(A\) est
\[
 E_A:=A+iA'.
\]
Ce n'est qu'en spécialisant \(A=\Xi\) et en supposant en outre que \(a\) est un zéro réel
simple de \(\Xi\) que l'on obtient
\[
 E_\Xi=\Xi+i\Xi'.
\]
\end{lemma}

\begin{proof}
On développe \(A\) et \(A'\) en séries de Taylor autour de \(a\), en utilisant
\(A(a)=0\) et \(A'(a)\ne0\), et l'on conserve les termes du premier ordre en
\(y\). Le résultat est strictement local : il ne démontre pas que \(E_A\) possède la
propriété globale de Hermite--Biehler.
\end{proof}

\begin{proposition}[Obstruction par les lois de comptage]
\label{pf84:E14}
Soit
\[
 N_A(T):=\Tr\mathbf1_{[-T,T]}(A).
\]
Les opérateurs \(Q\), \(B\), \(D_{\mathrm{orb}}\) et \(\mathcal R\) ne
possèdent pas à eux seuls la loi de comptage exigée d'un éventuel opérateur dual
des zéros.
\end{proposition}

\begin{proof}
Para \(T\ge1\),
\[
 N_Q(T)=\lfloor T\rfloor,
 \qquad
 N_B(T)=2\lfloor T\rfloor+1.
\]
Le théorème des nombres premiers donne
\[
 N_{D_{\mathrm{orb}}}(T)
 =2\sum_{p\le e^T}\left\lfloor\frac{T}{\log p}\right\rfloor
 \ge2\pi(e^T)\sim\frac{2e^T}{T}.
\]
Enfin, \(N_{\mathcal R}(T)\) est constante sur chaque intervalle
multiplicatif \([3^m,3^{m+1})\), car \(\mathcal R\) n'a que des niveaux
\(3^{m+1}\) de multiplicités finies.

Pour
\(\Xi(z)=\xi(1/2+iz)\), l'équation fonctionnelle implique
\(\Xi(-z)=\Xi(z)\) ; les zéros \(\pm\gamma\) ont la même multiplicité et
\[
 N_{\mathrm{sym}}(T)=2N_+(T)+O(1).
\]
Par conséquent, la convention bilatérale exige
\[
 N_{\mathrm{sym}}(T)\sim\frac{T}{\pi}\log T,
\]
tandis que la convention unilatérale équivalente, sous cette symétrie et cet appariement des
multiplicités, exige
\[
 N_+(T)\sim\frac{T}{2\pi}\log T.
\]
Aucune des quatre fonctions précédentes n'a cette asymptotique.
\end{proof}

\section{Laplaciens sur les cycles d'ordre \texorpdfstring{$3^m$}{puissances de trois}}

Pour $m\geq 1$, soient

\[
 N_m=3^m,
 \qquad
 C_m=\ZZ/N_m\ZZ,
\]

et considérons le laplacien combinatoire

\[
 (\Delta_m f)(j)=2f(j)-f(j+1)-f(j-1),
 \qquad f\in\ell^2(C_m).
\]

Les caractères

\[
 e_{m,n}(j)=N_m^{-1/2}
 \exp\!\left(\frac{2\pi i n j}{N_m}\right),
 \qquad n\in C_m,
\]

forment une base orthonormale et satisfont

\[
 \Delta_m e_{m,n}
 =4\sin^2\!\left(\frac{\pi n}{N_m}\right)e_{m,n}.
\]

Comme $N_m$ est impair, chaque classe de $C_m$ possède un unique représentant
$\widetilde n$ avec

\[
 -\frac{N_m-1}{2}\leq\widetilde n\leq\frac{N_m-1}{2}.
\]

Nous définissons l'opérateur autoadjoint $D_m$ par

\begin{equation}
 D_m e_{m,n}
 =\frac{N_m}{\sqrt\pi}
 \sin\!\left(\frac{\pi\widetilde n}{N_m}\right)e_{m,n}.
 \label{eq:operador-espectral-Dm}
\end{equation}

Alors

\begin{equation}
 D_m^2=\frac{N_m^2}{4\pi}\,\Delta_m.
 \label{eq:Dm-laplaciano}
\end{equation}

Soit $P_0$ la projection orthogonale sur le caractère constant. Pour toute
fonction paire bornée $F$ sur le spectre de $D_m$, nous introduisons la trace
positive

\begin{equation}
 \operatorname{Tr}^{+}_m\!\bigl(F(D_m)\bigr)
 :=\frac12\operatorname{Tr}
 \!\left(F(D_m)P_0^{\perp}\right).
 \label{eq:traza-positiva}
\end{equation}

Le facteur $1/2$ sélectionne une contribution par paire spectrale
$\{n,-n\}$ ; l'élimination de $P_0$ exclut la valeur propre nulle.

\begin{theorem}[Limite thêta des traces de chaleur]
\label{thm:limite-theta-triadico}
Pour tout $x>0$,
\[
 \lim_{m\to\infty}
 \operatorname{Tr}^{+}_m\!\left(e^{-xD_m^2}\right)
 =\Theta_+(x)
 :=\sum_{n\geq1}e^{-\pi n^2x}.
\]
\end{theorem}

\begin{proof}
La trace bilatérale, une fois le mode constant retiré, est
\[
 \sum_{\substack{n\in C_m\\n\neq0}}
 \exp\!\left[
 -\frac{N_m^2x}{\pi}
 \sin^2\!\left(\frac{\pi\widetilde n}{N_m}\right)
 \right].
\]
Pour chaque entier fixé $n$, le terme converge vers $e^{-\pi n^2x}$. De plus,
pour $|\widetilde n|\leq N_m/2$, l'inégalité
$\sin y\geq 2y/\pi$, valable sur $0\leq y\leq\pi/2$, fournit une
majorante gaussienne sommable et indépendante de $m$. La convergence
dominée appliquée à la somme discrète implique
\[
 \lim_{m\to\infty}
 \operatorname{Tr}\!\left(e^{-xD_m^2}P_0^\perp\right)
 =\sum_{n\in\ZZ\setminus\{0\}}e^{-\pi n^2x}.
\]
La division par deux établie dans \eqref{eq:traza-positiva} conserve une
contribution de chaque paire $\{n,-n\}$ et donne $\Theta_+(x)$.
\end{proof}

\section{Transformée thêta--Mellin}

Pour $\Re s>1$, nous définissons

\begin{equation}
 \mathcal Z_{3}(s)
 :=
 \frac{\displaystyle
 \int_0^\infty\Theta_+(x)x^{s/2-1}\,\dd x}
 {\pi^{-s/2}\Gamma(s/2)}.
 \label{eq:funcional-theta-mellin}
\end{equation}

La convergence absolue permet d'intégrer terme à terme. Comme

\[
 \int_0^\infty e^{-\pi n^2x}x^{s/2-1}\,\dd x
 =\Gamma(s/2)(\pi n^2)^{-s/2},
\]

on obtient

\begin{equation}
 \mathcal Z_3(s)=\sum_{n\geq1}n^{-s}=\zeta(s),
 \qquad \Re s>1.
 \label{eq:Z3-zeta}
\end{equation}

L'égalité se prolonge méromorphiquement par le prolongement de la fonction
zêta de Riemann \cite{Apostol1976}. Dans la suite, toutes les évaluations
proviennent de \eqref{eq:funcional-theta-mellin}, de ses dérivées ou de la
décomposition résiduelle de sa série de Dirichlet.

\begin{theorem}[Inversion thêta et équation fonctionnelle]
\label{pf84:C05}
Avec la notation opératorielle précédente,
\[
 \Theta_B(x):=\Tr(e^{-\pi xB^2})
 =1+2\Theta_+(x),
 \qquad
 Z(s):=\Tr(Q^{-s})=\mathcal Z_3(s)
 \quad(\Re s>1).
\]
La sommation de Poisson et la transformée de Mellin produisent
\[
 \Theta_B(x)=x^{-1/2}\Theta_B(1/x)
\]
et le prolongement de
\[
 \xi(s)=\frac12s(s-1)\pi^{-s/2}\Gamma(s/2)Z(s),
 \qquad
 \xi(s)=\xi(1-s).
\]
\end{theorem}

\begin{proof}
L'égalité \(\Theta_B=1+2\Theta_+\) sépare le mode zéro et apparie
\(\pm n\) ; l'identité \(Z=\mathcal Z_3\) est
\eqref{eq:Z3-zeta}. La sommation de Poisson appliquée à la gaussienne donne
l'inversion de \(\Theta_B\). Dans le demi-plan de convergence, l'intégration
terme à terme fournit
\[
 \pi^{-s/2}\Gamma(s/2)Z(s)
 =\int_0^\infty x^{s/2-1}\Theta_+(x)\,\dd x.
\]
En divisant l'intégrale en \((0,1)\) et \((1,\infty)\), en appliquant l'inversion
thêta à la première partie et en séparant les deux pôles classiques, on obtient le
prolongement et l'équation fonctionnelle. Cette unique preuve classique incorporée
ne localise pas les zéros de \(\zeta\).
\end{proof}

\section{Décomposition résiduelle modulo trois}

Pour $r\in\{0,1,2\}$, soit

\begin{equation}
 Z_r(s)=
 \sum_{\substack{n\geq1\\n\equiv r\pmod3}}n^{-s}.
 \label{eq:canales-residuales-Zr}
\end{equation}

En termes de la zêta de Hurwitz,

\[
 Z_0(s)=3^{-s}\zeta(s),
 \qquad
 Z_1(s)=3^{-s}\zeta\!\left(s,\frac13\right),
 \qquad
 Z_2(s)=3^{-s}\zeta\!\left(s,\frac23\right).
\]

Les trois termes ont pour résidu $1/3$ en $s=1$. Nous écrivons leurs
développements de Laurent sous la forme

\begin{equation}
 Z_r(1+z)-\frac1{3z}=\sum_{n\geq0}c_{r,n}\,z^n,
 \qquad
 \zeta(1+z)-\frac1z=\sum_{n\geq0}a_n\,z^n.
 \label{eq:coeficientes-Laurent-residuales}
\end{equation}

\begin{theorem}[Décomposition des coefficients de Laurent]
\label{thm:descomposicion-laurent-triadica}
Pour tout $n\geq0$,
\[
 a_n=c_{0,n}+c_{1,n}+c_{2,n},
 \qquad
 \gamma_n=(-1)^n n!\,a_n.
\]
Si $\chi_{-3}$ désigne le caractère réel non principal modulo trois, alors
\[
 \frac{L^{(n)}(1,\chi_{-3})}{n!}=c_{1,n}-c_{2,n}.
\]
De plus, avec \(\lambda=\log3\),
\[
 c_{0,n}=\frac13\left[
 \frac{(-\lambda)^{n+1}}{(n+1)!}
 +\sum_{k=0}^{n}a_k\frac{(-\lambda)^{n-k}}{(n-k)!}
 \right].
\]
Ces trois identités déterminent de manière unique
\((c_{0,n},c_{1,n},c_{2,n})\) à partir de
\((a_n,L^{(n)}(1,\chi_{-3}))\) et des coefficients agrégés d'ordre
inférieur.
\end{theorem}

\begin{proof}
La partition des entiers positifs en classes résiduelles fournit
$\zeta=Z_0+Z_1+Z_2$. En soustrayant les pôles en $s=1$ et en comparant les
coefficients de \eqref{eq:coeficientes-Laurent-residuales}, on obtient la
première identité. La deuxième coïncide avec la convention usuelle pour les
constantes de Stieltjes. De plus,
$Z_1-Z_2=L(s,\chi_{-3})$ ; les pôles de $Z_1$ et $Z_2$ s'annulent et la
comparaison des coefficients donne la troisième relation. Enfin,
\[
 Z_0(1+z)=\frac13e^{-\lambda z}\zeta(1+z).
\]
Le produit des deux séries de Laurent fournit la formule explicite
pour \(c_{0,n}\). La somme et la différence retrouvent alors
\(c_{1,n}\) et \(c_{2,n}\), ce qui prouve l'unicité affirmée.
\end{proof}

La somme des trois termes retrouve le secteur agrégé ; la différence des
classes $1$ et $2$ conserve l'orientation du caractère ; la comparaison
avec la classe divisible par trois enregistre le changement d'échelle. Ces trois
combinaisons sont linéairement indépendantes et déterminent les trois
coefficients $c_{r,n}$.

\section{Euler--Mascheroni et constantes de Stieltjes}

Soit

\[
 C_r:=\operatorname*{FP}_{s=1}Z_r(s).
\]

Les identités classiques pour la fonction digamma donnent
\cite{Apostol1976,Lagarias2013}

\begin{equation}
 C_0=\frac{\gamma-\log3}{3},
 \label{eq:C0-parte-finita}
\end{equation}

\begin{equation}
 C_1=\frac\gamma3+\frac{\log3}{6}+\frac{\pi}{6\sqrt3},
 \qquad
 C_2=\frac\gamma3+\frac{\log3}{6}-\frac{\pi}{6\sqrt3}.
 \label{eq:C1-C2-partes-finitas}
\end{equation}

\begin{corollary}
\label{cor:gamma-log3-L1}
Les parties finies satisfont
\[
 \boxed{\gamma=C_0+C_1+C_2},
\]
\[
 \boxed{L(1,\chi_{-3})=C_1-C_2=\frac{\pi}{3\sqrt3}},
\]
\[
 \boxed{\log3=C_1+C_2-2C_0}.
\]
\end{corollary}

\begin{proof}
À l'ordre zéro, le \cref{thm:descomposicion-laurent-triadica} donne
$a_0=c_{0,0}+c_{1,0}+c_{2,0}$. Par définition,
$a_0=\gamma$ et $c_{r,0}=C_r$, ce qui prouve la première égalité. La
différence des deux classes orientées est
$C_1-C_2=L(1,\chi_{-3})$. En substituant
\eqref{eq:C0-parte-finita}--\eqref{eq:C1-C2-partes-finitas}, on obtient les
deux autres identités.
\end{proof}

Une unique extraction des coefficients au moyen de la formule intégrale de
Cauchy appliquée à \eqref{eq:coeficientes-Laurent-residuales} détermine
simultanément $\gamma_0,\ldots,\gamma_6$. Les discrétisations à 96 et
192 nœuds coïncident jusqu'à l'ordre $10^{-87}$ ; l'identité des
coefficients provient toutefois du
\cref{thm:descomposicion-laurent-triadica} et non de cette comparaison numérique.

\section{Évaluation de la constante d'Apéry}

\begin{proposition}
\label{prop:evaluacion-apery}
La fonctionnelle \eqref{eq:funcional-theta-mellin} satisfait
\[
 \boxed{
 \mathcal Z_3(3)=\zeta(3)
 =1.202056903159594285399738161511449\ldots}.
\]
\end{proposition}

\begin{proof}
L'égalité exacte est la spécialisation $s=3$ de \eqref{eq:Z3-zeta}.
Pour certifier tout préfixe décimal sans modifier la fonctionnelle,
considérons $S_N=\sum_{n=1}^{N}n^{-3}$. Le critère intégral donne
\[
 S_N+\frac1{2(N+1)^2}
 <\zeta(3)<
 S_N+\frac1{2N^2}.
\]
Les intervalles rationnels emboîtés résultants déterminent les chiffres
imprimés.
\end{proof}

L'irrationalité de $\zeta(3)$ est le théorème classique d'Apéry
\cite{Apery1979} ; elle ne se déduit pas de la seule évaluation de la fonctionnelle.

\section{Constante de Glaisher--Kinkelin}

La dérivée spectrale en $s=-1$ définit

\begin{equation}
 \boxed{
 A_{\mathrm{GK}}
 =\exp\!\left(\frac1{12}-\mathcal Z_3'(-1)\right)}.
 \label{eq:glaisher-kinkelin-espectral}
\end{equation}

C'est la formule du déterminant régularisé associé au produit
hyperfactoriel \cite{Vardi1988}. Une dérivation centrale à cinq points,
évaluée à deux échelles, fournit

\[
 A_{\mathrm{GK}}
 =1.2824271291006226368753425688697917\ldots.
\]

Le terme $1/12$ appartient à la régularisation de l'hyperfactorielle et n'est pas
introduit comme paramètre d'ajustement.

\section{Constante d'Euler--Kronecker du corps d'Eisenstein}

Le caractère $\chi_{-3}$ détermine le corps quadratique imaginaire

\[
 F_{-3}=\QQ(\sqrt{-3}),
 \qquad
 \zeta_{F_{-3}}(s)=\zeta(s)L(s,\chi_{-3}).
\]

Sa constante d'Euler--Kronecker est

\begin{equation}
 \boxed{
 \gamma_{F_{-3}}
 =\gamma+\frac{L'(1,\chi_{-3})}{L(1,\chi_{-3})}}.
 \label{eq:euler-kronecker-eisenstein}
\end{equation}

Les coefficients résiduels de
\eqref{eq:coeficientes-Laurent-residuales} donnent

\[
 L(1,\chi_{-3})
 =0.6045997880780726168646927525473852\ldots,
\]
\[
 L'(1,\chi_{-3})
 =0.2226629869686015094866602627647443\ldots,
\]

et, par \eqref{eq:euler-kronecker-eisenstein},

\[
 \gamma_{F_{-3}}
 =0.9454972808716807032397499941581890\ldots.
\]

Le choix de $F_{-3}$ provient du caractère résiduel modulo trois, et non
d'une recherche sur les valeurs décimales de constantes de corps
quadratiques.

\section{Produit local associé aux nombres premiers jumeaux}

Pour un motif fini $H=\{h_1,\ldots,h_k\}\subset\ZZ$, nous définissons

\[
 \nu_p(H):=\left|\{-h_1,\ldots,-h_k\}\bmod p\right|,
 \qquad
 \sigma_p(H)=
 \frac{1-\nu_p(H)/p}{(1-1/p)^k}.
\]

Pour $H=\{0,2\}$, les facteurs locaux conduisent à

\begin{equation}
 C_2^{\mathrm{twin}}
 =\prod_{p>2}\frac{p(p-2)}{(p-1)^2}.
 \label{eq:constante-primos-gemelos}
\end{equation}

Soit \(P_X\) le produit restreint aux nombres premiers \(2<p\leq X\). Pour
\(p>X\), tous les facteurs appartiennent à \((0,1)\) ; en étendant le produit
de la queue des nombres premiers à tous les entiers, on obtient la borne
télescopique
\[
 \prod_{n>X}\left(1-\frac1{(n-1)^2}\right)
 =\frac{X-1}{X}
 \leq
 \prod_{p>X}\left(1-\frac1{(p-1)^2}\right)<1.
\]
Le produit \(P_{2\times10^6}\), évalué avec arrondi décimal dirigé, et
cette inégalité rationnelle pour la queue donnent

\begin{equation}
 0.6601615071224244826403933551
 <C_2^{\mathrm{twin}}<
 0.6601618372035081248537566591.
 \label{eq:intervalo-primos-gemelos}
\end{equation}

Le décompte direct jusqu'à $10^6$ contient 8169 paires de nombres premiers jumeaux. La
convergence du produit local et ce recensement fini n'impliquent pas
l'existence d'une infinité de paires ; cette dernière est une proposition arithmétique
distincte.

\section{Constante de Meissel--Mertens et hamiltonien premier}

La constante de Meissel--Mertens est définie par

\begin{equation}
 B_1=\lim_{x\to\infty}
 \left(\sum_{p\leq x}\frac1p-\log\log x\right)
 \label{eq:meissel-mertens-def}
\end{equation}

et satisfait

\begin{equation}
 B_1=\gamma+
 \sum_p\left(\log\!\left(1-\frac1p\right)+\frac1p\right).
 \label{eq:meissel-mertens-euler}
\end{equation}

Avec la notation de \cref{pf84:C08}, nous posons
\[
 \mathcal F_P:=\Gamma_s(\mathfrak h_{\mathbb P}),
 \qquad
 H_P:=d\Gamma(H_{\mathbb P}).
\]
Ni l'espace de Fock bosonique des modes premiers ni son hamiltonien ne sont redéfinis ici.
L'opérateur déjà construit a pour fonction de partition

\begin{equation}
 \operatorname{Tr}_{\mathcal F_P}(e^{-sH_P})
 =\prod_p\frac1{1-p^{-s}}
 =\zeta(s),
 \qquad \Re s>1.
 \label{eq:hamiltoniano-primo-zeta}
\end{equation}

Soit \(\mathcal H_P^{(1)}\) le sous-espace à une particule, engendré par
\(\{|p\rangle:p\text{ premier}\}\). Alors
\[
 \operatorname{Tr}_{\mathcal H_P^{(1)}}\!\left(
 e^{-H_P}\mathbf 1_{(-\infty,\log x]}(H_P)\right)
 =\sum_{p\leq x}\frac1p.
\]
Ainsi, \(B_1\) est la partie finie, lorsque \(x\to\infty\), de cette trace
spectrale tronquée après soustraction de \(\log\log x\). L'égalité
\eqref{eq:hamiltoniano-primo-zeta} présuppose la décomposition de Fock sur
les nombres premiers ; elle constitue donc une représentation opératorielle du produit
d'Euler, non une déduction de la primalité à partir du laplacien
triadique.

\section{Dépendance à l'égard de l'opérateur}

Les évaluations précédentes partagent l'opérateur
\eqref{eq:operador-espectral-Dm}, la trace positive
\eqref{eq:traza-positiva} et la transformée
\eqref{eq:funcional-theta-mellin}. D'autres constantes exigent des dynamiques
distinctes. La constante de Catalan est associée au caractère modulo quatre
$\chi_{-4}$, tandis que la constante de Khinchin appartient à l'opérateur de
Gauss et à sa mesure invariante. Un préfixe de l'une quelconque d'entre elles peut
apparaître dans un catalogue fini sans que la fonctionnelle triadique la sélectionne.
Cette dépendance opératorielle distingue l'évaluation spectrale de la
capacité purement combinatoire de codage.
